% !TEX root = thesis_defense_0421052099.tex
% Frames only -- no \documentclass and no \begin{document}. Build the root
% named above; this file on its own stops with "Emergency stop".
%
% Prose rule, the same one the book follows since 2026-09-30: no em-dashes
% and no colon explanations. A thought that needs a dash is two sentences.
% Colons stay where they introduce a list, and after a bold or alert label.

% ===================== 1 =====================
{\setbeamertemplate{footline}{}
\begin{frame}[plain]
  \titlepage
\end{frame}}

% ===================== 2 =====================
% Slides 2 and 3 of the August deck, merged. The definition and the two
% kinds of hallucination are the whole of what the rest of the talk needs
% from this part; the separate "the problem" slide said the first sentence
% of the definition twice.
\begin{frame}{The problem}
  Language models answer just as fluently when they \alert{do not hold the
  fact}.

  A \emph{hallucination} is output with \alert{no source the model can
  point to}.
  \vspace{2mm}
  \begin{columns}[T]
    \begin{column}{0.47\textwidth}
      \begin{block}{Factually wrong}
        \agrraggedright
        It contradicts the world (\emph{``Ada Lovelace was born in
        Paris''}). To catch it you need the true answer already, which is
        the thing we are trying to produce.
      \end{block}
    \end{column}
    \begin{column}{0.47\textwidth}
      \begin{block}{Unsupported}
        \agrraggedright
        It may even be true, but nothing the system retrieved says so. To
        catch it you need \alert{only what was retrieved}, and we have that.
      \end{block}
    \end{column}
  \end{columns}
  \takeaway{A retrieval system can only police the second kind. That is the
    one this thesis attacks, and the only one it claims.}
\end{frame}

% ===================== 3 =====================
% Slides 4 and 5 merged: the three origins, and which of them a deployer can
% act on. "And multi-hop compounds it" moved to slide 4, where the knowledge
% graph slide used to make the same point a second time.
\begin{frame}{Where it comes from, and what we can change}
  Not a bug that escaped testing. It comes from how these models are built,
  so the response has to be architectural.
  \vspace{3mm}
  \begin{columns}[T]
    \begin{column}{0.31\textwidth}
      \begin{block}{The training data}
        \agrraggedright
        Errors and thin coverage, reproduced faithfully. Fixing them means
        retraining.
      \end{block}
    \end{column}
    \begin{column}{0.33\textwidth}
      \begin{block}{The model}
        \agrraggedright
        Knowledge is spread across the weights with \mbox{\alert{no index}}.
        Nothing separates a memorised fact from a plausible continuation.
      \end{block}
    \end{column}
    \begin{column}{0.31\textwidth}
      \begin{block}{\alert{The prompt}}
        \agrraggedright
        What enters the context window. It is the one origin a deployer
        controls.
      \end{block}
    \end{column}
  \end{columns}
  \takeaway{Control what goes in, then check what comes out against it,
    before anything is emitted.}
\end{frame}

% ===================== 4 =====================
% Slides 6 and 7 merged. The mediator block keeps its one number, which
% environment.tex's graph-statistics table gives as 1,652,618 nodes (63.7%).
\begin{frame}{Knowledge graphs, and why multi-hop is hard}
  \begin{columns}[T]
    \begin{column}{0.52\textwidth}
      \agrraggedright
      Facts are stored as \alert{triples} of head, relation and tail.
      \smallskip

      \hspace{3mm}\textsf{(Ada Lovelace, place\_of\_birth, London)}
      \smallskip

      Every fact has an address, so a claim can be checked by looking for
      an edge.
      \smallskip

      A question that chains two facts is \alert{multi-hop}. Each hop
      starts from the last, so a wrong first hop is never revisited.
    \end{column}
    \begin{column}{0.44\textwidth}
      \begin{block}{The complication: mediators}
        \agrraggedright
        Freebase stores a fact with more than two participants as an
        unnamed \alert{node}. ``Rainn Wilson played Dwight in
        \mbox{\textsf{The Office}}'' is one. $63.7\%$ of this graph's nodes
        are mediators, so a path that reads as two hops is often three.
      \end{block}
    \end{column}
  \end{columns}
  \takeaway{Nothing in the graph is called \textsf{plays}. The relation a
    question names is rarely the relation the graph stores.}
\end{frame}

% ===================== 5 =====================
\begin{frame}{Where existing approaches stop}
  % Widths measured from the cells, not guessed. The last column cannot be
  % made single-line -- "Emits whatever the last generation call produces" is
  % 78mm on its own -- but at 45mm no row orphans one word onto a second line.
  % Rules between the rows rather than wider columns: three of the four
  % "Where it stops" cells run to two lines, and widening cannot remove the
  % wrap inside 140mm, so the fix is the row boundary.
  \begin{table}
    \small
    \setlength{\tabcolsep}{4pt}
    % 1.1 rather than 1.2 pays for the precision sentence under the table.
    \renewcommand{\arraystretch}{1.1}
    \begin{tabular}{@{}L{31mm}L{56mm}L{45mm}@{}}
      \toprule
      \textbf{Approach} & \textbf{What it does} & \textbf{Where it stops} \\
      \midrule
      Parametric        & Answers from weights
                        & No external evidence at all \\
      \midrule
      Vector-RAG        & Retrieves text once, then reasons
                        & Retrieval happens \emph{before} reasoning begins \\
      \midrule
      Static GraphRAG   & Retrieves a fixed neighbourhood
                        & Radius fixed before the question is understood \\
      \midrule
      Agentic navigation & Interleaves retrieval and reasoning
                        & Emits whatever the last generation call produces \\
      \bottomrule
    \end{tabular}
  \end{table}
  % sec:limitations: "The problem is not fabrication. The problem is
  % precision." Said here, slide 20's zero for both navigating systems is
  % the expected result rather than a surprise that undercuts the opening.
  {\small Agentic systems already name only what they traversed. Their
  problem is \alert{precision}. A real, traversed entity can still be the
  wrong answer.\par}
  \takeaway{None of the four checks what the answer asserts \emph{before}
    it is delivered.}
\end{frame}

% ===================== 6 =====================
% "Beam search" arrives on the next slide as Think-on-Graph's row, and it
% has to be defined before then. The takeaway plants the limitation that
% slide 9's Backtracker answers, and "multiplies" is what slide 19 cashes.
\begin{frame}{What guided beam search is}
  \begin{columns}[T]
    \begin{column}{0.56\textwidth}
      \small
      It keeps only the most promising branches:
      \begin{enumerate}\itemsep2pt
        \item List every edge out of the current \alert{frontier}.
        \item \alert{Score} each for how likely it is to help.
        \item Keep the best $b$, the \alert{beam width}.
        \item Expand those, and repeat to a depth cap.
      \end{enumerate}
      \vspace{1mm}
      \emph{Guided} means the score is a judgement.
    \end{column}
    \begin{column}{0.40\textwidth}
      \begin{block}{Why bound it}
        \small\agrraggedright
        Entities carry hundreds of relations, so three unbounded hops is
        millions of paths. When the scorer is a \alert{model call}, the
        width \emph{is} the cost, and it multiplies with every hop.
      \end{block}
    \end{column}
  \end{columns}
  \takeaway{Everything outside the beam is discarded unseen, and a plain
    beam has no way back for it.}
\end{frame}

% ===================== 7 =====================
\begin{frame}{What everyone else does}
  \bodyshift{-2.75pt}
  \begin{table}
    \small
    \setlength{\tabcolsep}{5pt}
    % Plan-on-Graph in full. Paths-over-Graph is also cited as "PoG" and
    % reports higher numbers, and background.tex writes both names out for
    % that reason. AGR is not a row: it is introduced on slide 9.
    \begin{tabular}{@{}lL{42mm}ccc@{}}
      \toprule
      \textbf{System} & \textbf{Exploration} & \textbf{Decomp.} &
      \textbf{Backtrack} & \textbf{Verifies first} \\
      \midrule
      GraphRAG       & Fixed radius, no search  & No  & No  & No \\
      RoG            & Generated relation paths & Partial & No & No \\
      Think-on-Graph & LLM-guided beam search   & No  & No  & No \\
      Plan-on-Graph  & Adaptive planning        & Yes & Yes & Partial \\
      \bottomrule
    \end{tabular}
  \end{table}
  % The two qualifications background.tex attaches to this table. RoG's
  % answers are faithful by construction (tab:mechanisms), and the
  % post-hoc checkers are where the decomposition step comes from
  % (sec:verify-motivation), so naming them is cheaper than being asked.
  {\small\agrraggedright RoG's answers are grounded by construction, but
  nothing examines them. Chain-of-Verification and FActScore check claims
  \emph{after} generation. Published Hits@1 is higher and \emph{not}
  comparable, because of other backbones, subsets and full Freebase.\par}
  \takeaway{No surveyed graph-navigating system checks its answer's claims
    against the triples it retrieved, before answering.}
\end{frame}

% ===================== 8 =====================
% The questions as sec:rqs words them, word for word, and the six
% contributions as a preview. Slide 29 is the same six, in the same order.
% RQ1 used to carry "and does the advantage grow with hop count", which is
% sec:rqs's elaboration rather than its question; it is slide 18's subtitle
% now. check_slides.py holds all three to the book's bold lines.
\begin{frame}{Research questions}
  \bodyshift{2.75pt}
  % Ragged inside both lists. Beamer's list code re-issues \raggedright on
  % entry, which the preamble repoints at \justifying, so a list in a narrow
  % column is justified unless \raggedright is put back for the column. At
  % this measure justified items hyphenated "contribute" and spaced
  % "AGR and its" across the whole line.
  \begin{columns}[T]
    \begin{column}{0.58\textwidth}
      \let\raggedright\agrraggedright
      % A margin as wide as the label, so the labels start at the column's
      % edge. At the default margin they hung 6pt outside it, and the body
      % sat 4.7pt left of centre.
      \settowidth{\leftmargini}{\textbf{RQ1}}%
      \addtolength{\leftmargini}{\labelsep}%
      \begin{enumerate}
        \item[\textbf{RQ1}] Does agentic navigation improve
          \alert{multi-hop} factual accuracy?
          \vspace{1mm}
        \item[\textbf{RQ2}] What does pre-generation verification
          \alert{contribute} beyond graph navigation?
          \vspace{1mm}
        \item[\textbf{RQ3}] Which components contribute what, at what
          \alert{token cost}?
      \end{enumerate}
      \vspace{1mm}
      {\small\agrraggedright RQ2 asks what verification
      \emph{contributes}, not whether it reduces hallucination. The answer
      has several parts.\par}
    \end{column}
    \begin{column}{0.38\textwidth}
      \let\raggedright\agrraggedright
      \begin{block}{Six contributions}
        \small
        \begin{enumerate}\setlength{\itemsep}{0pt}
          \item AGR and its verification layer
          \item A component-level ablation
          \item Stratum-dependent decomposition
          \item The echo attractor
          \item Benchmark-defect rates
          \item A pre-specified protocol
        \end{enumerate}
      \end{block}
    \end{column}
  \end{columns}
  \centrebody
\end{frame}

% ===================== 9 =====================
\begin{frame}{AGR: An explicit state machine}
  \begin{columns}[T]
    \begin{column}{0.46\textwidth}
      \small\let\raggedright\agrraggedright\agrraggedright
      Six nodes. The model is consulted only at decision points.
      % The Backtracker restores the HIGHEST-scoring snapshot, "not the most
      % recent one, which would make backtracking a simple undo"
      % (sec:backtracking), and bans what the last pass took. The Verifier
      % is one part of the architecture, not "the contribution" on its own
      % (the opening of ch:framework).
      % \parskip is 4pt deck-wide and applies between items too; at six
      % items that alone pushed the cycles sentence off the slide.
      \begin{itemize}\setlength{\itemsep}{0pt}\setlength{\parskip}{0pt}
        \item \textbf{Planner}: sub-objectives
        \item \textbf{Explorer}: best-first, beam 3
        \item \textbf{Evaluator}: objective met?
        \item \textbf{Backtracker}: best earlier frontier, ban list
        \item \textbf{Verifier}: checks claims first
        \item \textbf{Answerer}: emits what survived
      \end{itemize}
      % "Three cycles", counted from the arrows returning to the Explorer
      % in the tikzpicture below, and named after their edge labels so that
      % counting confirms the sentence. check_slides.py holds both.
      Three cycles (\emph{continue}, \emph{backtrack}, \emph{retry}), all
      bounded by budgets rather than by model behaviour.
    \end{column}
    \begin{column}{0.52\textwidth}
      \centering
      % Every gap here is load-bearing. An edge label carries fill=white, so a
      % gap shorter than the label hides the arrow it sits on. node distance
      % is set from the label heights, the backtrack run from the width of
      % the word "backtrack", and the two left-hand loops sit at different
      % depths with their labels above and below their own runs.
      %
      % Scale 0.70 with footnotesize labels, up from 0.60 and scriptsize: at
      % the old size the edge labels printed at 4.8pt. Uniform scaling keeps
      % every label-to-gap ratio, and narrower boxes (25mm, down from 27)
      % with shorter runs (17mm, 21mm and 24mm, down from 18, 26 and 28) pay
      % for the larger scale inside this column. Measured on the page: at an
      % 18mm run the "backtrack" label sat on the dashed border, and at 15mm
      % "continue" did.
      \begin{tikzpicture}[
        node distance=9mm, scale=0.70, transform shape,
        box/.style={draw=agrNode, thick, rounded corners=2pt,
          minimum height=6mm, fill=agrNode!10,
          text width=25mm, align=center, font=\small},
        vbox/.style={box, draw=agrSup, fill=agrSup!18},
        term/.style={draw=black!45, circle, inner sep=1.5pt,
          font=\footnotesize, fill=black!6},
        flow/.style={-{Stealth[length=2mm]}, semithick, draw=black!68},
        lbl/.style={font=\footnotesize, inner sep=1.2pt, fill=white}]
        \node[term]                    (start) {START};
        \node[box, below=6mm of start] (plan)  {Planner};
        \node[box, below=of plan]      (expl)  {Explorer};
        \node[box, below=of expl]      (eval)  {Evaluator};
        \node[vbox, below=of eval]     (ver)   {Verifier};
        \node[box, below=of ver]       (ans)   {Answerer};
        \node[box, right=21mm of eval] (back)  {Backtracker};
        \draw[flow] (start) -- (plan);
        \draw[flow] (plan)  -- (expl);
        \draw[flow] (expl)  -- (eval);
        \draw[flow] (eval)  -- node[lbl] {answer} (ver);
        \draw[flow] (ver)   -- node[lbl] {grounded $\mid$ give\_up} (ans);
        \draw[flow] (eval.east) -- node[lbl] {backtrack} (back.west);
        \draw[flow] (back.north) |- (expl.east);
        % search loop: label rides above its own run, clear of both the
        % Evaluator and the vertical it turns into.
        \draw[flow] (eval.west) -- ++(-17mm,0)
          node[lbl, above, pos=0.5] {continue} |- (expl.west);
        % re-exploration loop: deeper, label below, so the two never meet.
        \draw[flow] (ver.west) -- ++(-24mm,0)
          node[lbl, below, pos=0.5] {retry} |- ([yshift=-1.8mm] expl.west);
        \begin{scope}[on background layer]
          \node[draw=black!35, dashed, rounded corners, inner sep=1.2mm,
                fit=(expl)(eval)] {};
        \end{scope}
      \end{tikzpicture}
    \end{column}
  \end{columns}
  \centrebody
\end{frame}

% ===================== 10 =====================
\begin{frame}{Constrained tools, not free-form queries}
  % One agent, not six. The Node column says which of its six nodes issues
  % each call, and check_slides.py holds that column against the call sites
  % in agr/, so it cannot drift from the code.
  One agent with six nodes, and only three of them reach the graph, each
  through a fixed operation.
  % Natural column widths, not fixed ones: every cell here is a short phrase,
  % and four l columns cannot wrap, so the rows are single-line by
  % construction. The table comes to ~134mm of the 140mm text block.
  \begin{table}
    \small
    \begin{tabular}{@{}llll@{}}
      \toprule
      % "At most" says how much comes back, in the same form for every row.
      % Two of the four name no cap -- verify_connection returns a boolean
      % and search_entity runs a resolver -- so the header claims none.
      \textbf{Operation} & \textbf{Node} & \textbf{Returns} &
      \textbf{At most} \\
      \midrule
      % Run order, so the column reads as the path through the machine.
      % Every name in this column is checked against a def in kg_tools.py.
      \texttt{search\_entity}  & Planner  & Surface form $\rightarrow$ node
                               & $3$ candidates \\
      \texttt{get\_relations}  & Explorer & Candidate relations
                               & $\leq 300$ to the scorer \\
      \texttt{get\_neighbors}  & Explorer & Neighbours of an entity
                               & $\leq 200$ per expansion \\
      \texttt{verify\_connection} & Verifier & Adjacency, full graph
                               & one boolean \\
      \bottomrule
    \end{tabular}
  \end{table}
  {\small Evaluator, Backtracker and Answerer read state only. The agent
  never writes Cypher, and every call is logged with its arguments and
  result.\par}
  \takeaway{Determinism in the tools is what makes the verification layer
    possible.}
\end{frame}

% ===================== 11 =====================
% ch:framework opens by introducing the layer "as one part of the
% architecture rather than as the part the rest exists to serve", and says
% its measured contribution is "auditability rather than accuracy". The
% takeaway says both, so slides 20 to 22 read as the measurement this slide
% announced rather than as a retraction of it.
\begin{frame}{The Structural Verification Layer}
  \begin{columns}[T]
    \begin{column}{0.58\textwidth}
      \let\raggedright\agrraggedright
      Before anything is emitted:
      \begin{enumerate}
        \item The draft is split into \alert{atomic claims}
        \item Each claim is checked against the graph, the \alert{triples
          actually traversed} first
        \item Claims grounded by the traversal keep those triples as
          \alert{evidence}
        \item Unsupported claims send the agent back to
          \alert{re-explore}, or are \alert{dropped}
      \end{enumerate}
    \end{column}
    \begin{column}{0.39\textwidth}
      % The abstract's condition, kept: "where it cannot ground a claim,
      % AGR hedges rather than asserts". Unconditioned, the sentence says
      % the system always hedges.
      \begin{block}{The behaviour this buys}
        \agrraggedright
        Where it cannot ground a claim, the system \textbf{hedges rather
        than asserts}. Where the walk grounds it, the answer arrives with
        that evidence.
      \end{block}
    \end{column}
  \end{columns}
  \takeaway{The layer this thesis set out to contribute. Measured, it buys
    auditability rather than accuracy.}
\end{frame}

% ===================== 12 =====================
% The bounds are the supervisor's first two issues, and check_slides.py
% holds each of them on this slide AND in the script. They have been
% proposed for deletion once already, on the ground that the slide is wordy;
% the wordiness was never the bounds.
\begin{frame}{What \emph{structural} means, and what it does not}
  \bodyshift{-2pt}
  % "Against the graph the agent walked, not against the model's judgement"
  % was true of route 1 only. Route 2 asks the FULL graph, and route 3 is a
  % model call. framework.tex states the actual rule: the model "is
  % consulted only where the graph is silent ... It never overrides a
  % finding of the environment."
  The graph decides. The model is consulted only where the graph is
  silent, and it never overrides the graph.
  \vspace{2mm}
  \begin{columns}[T]
    \begin{column}{0.48\textwidth}
      \begin{block}{What it does not mean}
        It tests \alert{adjacency}, not the relation the claim names, and
        not its direction. A \emph{mother} claim passes on a \emph{child}
        edge.
      \end{block}
      \vspace{2mm}
      \begin{block}{What it cannot catch}
        \agrraggedright
        A claim can be \emph{true} and still be the wrong answer. That is
        the echo attractor of slide 27.
      \end{block}
    \end{column}
    \begin{column}{0.48\textwidth}
      \begin{block}{What ``its evidence'' means}
        Only the check against the \alert{walked} graph attaches triples.
        The other two routes accept with nothing attached. The log keeps
        the \emph{count}, not the triples.
      \end{block}
    \end{column}
  \end{columns}
  \centrebody
\end{frame}

% ===================== 13 =====================
% Drawn here rather than \input from thesis_book/figures/fig_claim_path.tex.
% That figure is a tall vertical chain, and a 16:9 slide is height-bound:
% scaled to fit, its labels printed at 5.2pt. This is the same routing laid
% out left to right, at the slide's own size, with the book figure's test
% wording verbatim -- check_slides.py holds every test label here against
% fig_claim_path.tex, so the two cannot drift apart. The exit rule the book
% figure draws as a box is the sentence under the picture.
\begin{frame}{One claim, three routes}
  \bodyshift{-1.25pt}
  \centering
  \vspace{2mm}
  % Sizes are the slide's own: footnotesize throughout, unscaled, so the
  % smallest label prints at 10pt. The verify_connection test is 4mm wider
  % than the others because its first line is one unbreakable token.
  \begin{tikzpicture}[
      node distance=7mm,
      execute at begin picture={\hyphenpenalty=10000\exhyphenpenalty=10000},
      box/.style={draw=agrNode, thick, rounded corners=2pt,
        minimum height=11mm, text width=24mm, align=center,
        font=\footnotesize, fill=agrNode!8},
      dec/.style={draw=agrNode!55, rounded corners=2pt, minimum height=13mm,
        text width=24mm, align=center, font=\footnotesize, fill=agrNode!4},
      bucket/.style={draw, rounded corners=5pt, minimum height=6mm,
        align=center, font=\footnotesize\sffamily},
      sup/.style={bucket, draw=agrSup, thick, fill=agrSup!14},
      uns/.style={bucket, draw=agrUns, thick, fill=agrUns!14,
        text width=22mm},
      flow/.style={-{Stealth[length=2mm]}, semithick, draw=black!68},
      lbl/.style={font=\footnotesize, inner sep=1.2pt, fill=white},
      ev/.style={font=\footnotesize\itshape, inner sep=1.2pt,
        text=agrSup}]
    % The main line: one claim, left to right, leaving at the first test
    % that settles it.
    \node[dec]                    (d1)  {both endpoint names seen
                                         during traversal?};
    \node[dec, right=7mm of d1]   (d2)  {pair in \emph{traversed}
                                         adjacency?};
    \node[dec, right=7mm of d2, text width=28mm]
                                  (d3)  {\textsf{verify\_connection}:
                                         adjacent in the full graph?};
    \node[box, right=7mm of d3]   (ent) {entailment check\\
                                         {\itshape\color{black!60}one
                                         batched call}};
    \node[box, above=7mm of d1]   (draft) {Draft and decompose\\
                                         {\itshape\color{black!60}one
                                         model call}};
    % Supported spans from the first structural test to the entailment
    % check, so each "yes" is a short vertical into it.
    \path let \p1 = (d2.west), \p2 = (ent.east), \p3 = (draft.center) in
      node[sup, minimum width=\x2-\x1, anchor=west] (sup) at (\x1,\y3)
      {supported $S$};
    % 12mm, not less: the skip lane below crosses at d1.south - 8mm, and
    % d1 is the deepest node, so a nearer bucket sits on the lane.
    \node[uns, below=12mm of ent] (uns) {unsupported $U$};

    \draw[flow] (draft) -- node[lbl] {claims $K$} (d1);
    \draw[flow] (d1) -- node[lbl] {yes} (d2);
    \draw[flow] (d2) -- node[lbl] {no} (d3);
    \draw[flow] (d3) -- node[lbl] {no} (ent);
    % The one route that records evidence.
    \draw[flow] (d2.north) -- node[lbl, pos=0.35] {yes}
      node[ev, right, pos=0.70] {+ evidence} (d2.north |- sup.south);
    \draw[flow] (d3.north) -- node[lbl, pos=0.35] {yes}
      (d3.north |- sup.south);
    \draw[flow] (ent.north) -- node[lbl, pos=0.35] {entailed}
      (ent.north |- sup.south);
    \draw[flow] (ent) -- node[lbl, right] {not entailed} (uns);
    % Endpoint never traversed: both structural routes are skipped. The
    % lane runs under the dashed region and its caption, and enters the
    % entailment check left of the arrow that leaves it.
    \draw[flow] (d1.south) -- ++(0,-8mm) node[lbl, pos=0.5] {no}
      -| ([xshift=-7mm]ent.south);
    \begin{scope}[on background layer]
      \node[draw=black!35, dashed, rounded corners, inner sep=2mm,
            fit=(d2)(d3)] (strbox) {};
    \end{scope}
    \node[font=\footnotesize\itshape, text=black!70, below=0.6mm of strbox]
      {structural routes};
  \end{tikzpicture}

  \vspace{2mm}
  {\small Only the third route costs a model call. Neither lookup reads the
  relation. The node exits \textsf{grounded} if $U$ is empty, else
  \textsf{retry} while budget remains, else \textsf{give\_up}.\par}
  \centrebody
\end{frame}

% ===================== 14 =====================
% WebQTest-960, read from results/phase4/test_webqsp_agr.jsonl, and
% check_slides.py holds every figure on this slide to that record. The
% mention "The Office" is the dataset's annotation (use_gold_entities); the
% planner's search_entity call resolves it to a node. Each Explorer pass
% keeps three relations, and the one named is the highest-scoring. The
% mention is quoted because it is a string until it resolves: as
% \textsf{The Office} it read "the given mention The Office", since \textsf
% is invisible in a deck already set in sans.
\begin{frame}{One question, end to end}
  \bodyshift{-0.75pt}
  \small
  \textsf{WebQTest-960}, ``who plays dwight in the office''
  \vspace{1mm}
  \begin{columns}[T]
    \begin{column}{0.60\textwidth}
      \let\raggedright\agrraggedright
      \begin{enumerate}\setlength{\itemsep}{0pt}
        \item \alert{Planner}: two sub-objectives. The given mention
          ``The Office'' resolves to one node.
        \item \alert{Explorer}: keeps three relations. The best is\\
          \textsf{tv.tv\_program.regular\_cast}
        \item \alert{Evaluator}: objective met, resolves
          \textsf{Dwight Schrute}
        \item \alert{Explorer}: keeps three again. The best is\\
          \textsf{tv.regular\_tv\_appearance.character}
        \item \alert{Evaluator}: resolves \textsf{Rainn Wilson}, decides
          \textsf{answer}
        \item \alert{Verifier}: $2$ claims, $2$ supported, none dropped
      \end{enumerate}
    \end{column}
    \begin{column}{0.38\textwidth}
      % n_supporting_triples counts matches, not distinct triples, so it is
      % an upper bound (sec:output-contract). "At most" is that bound.
      \begin{block}{What came out}
        \small\agrraggedright
        ``Rainn Wilson plays Dwight Schrute in The Office.''
        \smallskip

        $6$ model calls \textperiodcentered\ depth $2$ \textperiodcentered\
        at most $18$ supporting triples
      \end{block}
      \vspace{1mm}
      {\small\agrraggedright \alert{Both} hops run \emph{through} a
      mediator. The relation named is the edge \emph{into} the appearance
      record, not the one that lands on the entity.\par}
    \end{column}
  \end{columns}
  \centrebody
\end{frame}

% ===================== 15 =====================
% Two disclosures the old slide lacked, both from the book. The graph is a
% union of subgraphs extracted to contain their answers, which threat 1 in
% sec:threats says "bounds every absolute number". And reachability
% certifies entity answers only: sec:unanswerable defines a class of
% questions the graph cannot answer at all, so "a miss is the system's,
% not the environment's" was not true.
\begin{frame}{The environment and the question sets}
  \bodyshift{-2.75pt}
  \begin{columns}[T]
    \begin{column}{0.45\textwidth}
      % Ruled and headed like the table beside it, so the block title does
      % not read as a header row.
      %
      % Both tables run the full width of their column. At their natural
      % widths they sat at the columns' left edges, 45pt and 33pt short of
      % the right ones, and the slide's ink ran from 25pt off the left edge
      % of the page to 60pt off the right.
      \begin{block}{Freebase-derived graph}
        \small
        \begin{tabular*}{\linewidth}{@{\extracolsep{\fill}}lr@{}}
          \toprule
          \textbf{Property}  & \textbf{Value} \\
          \midrule
          Entities           & $2{,}592{,}892$ \\
          Triples            & $8{,}309{,}194$ \\
          Distinct relations & $7{,}058$ \\
          Import time        & $36.4$\,s \\
          \bottomrule
        \end{tabular*}
      \end{block}
      {\small\agrraggedright A union of RoG's per-question subgraphs, each
      cut to contain its answer, so friendlier than full Freebase.\par}
    \end{column}
    \begin{column}{0.5\textwidth}
      \begin{block}{Certified question sets}
        \small
        \begin{tabular*}{\linewidth}{@{\extracolsep{\fill}}lrr@{}}
          \toprule
                       & \textbf{WebQSP} & \textbf{CWQ} \\
          \midrule
          Questions      & $400$ & $400$ \\
          Gold (median)  & $1.5$ & $1.0$ \\
          Reachable      & $97.0\%$ & $99.2\%$ \\
          Multi-hop      & $132$ & $260$ \\
          \bottomrule
        \end{tabular*}
      \end{block}
      {\small\agrraggedright Reachability caps every Hits@1, for entity
      answers only. Dates, quantities and superlatives are beyond this
      graph.\par}
    \end{column}
  \end{columns}
  \takeaway{WebQSP and ComplexWebQuestions (CWQ), unmodified. Absolute
    levels are optimistic, and every system faces the same graph.}
\end{frame}

% ===================== 16 =====================
% sec:baseline-constants names TWO things not held constant: the candidate
% widths, and the drafting prompt, whose grounding provisions the baselines'
% plain prompt lacks. The takeaway used to call the widths "the one
% exception". It also claimed the differences were not due to "spend",
% which slide 17 contradicts: AGR spends more tokens than Think-on-Graph.
%
% What must not go: "Narrower is cheaper" and "lower bound". Disclosing that
% the widths differ, without saying which way the difference cuts, invites
% the inference that the confound could explain ToG's clipping. It cannot.
% check_slides.py holds both phrases and both pairs of numbers, read from
% agr/baselines/tog.py and agr/kg_tools.py.
\begin{frame}{Making the comparison fair}
  \begin{columns}[T]
    \begin{column}{0.44\textwidth}
      \small\let\raggedright\agrraggedright\agrraggedright
      Five systems, \alert{one frozen backbone}, \alert{one budget}:
      \begin{itemize}\setlength{\itemsep}{0pt}\setlength{\parskip}{0pt}
        \item No-retrieval (parametric control)
        \item Vector-RAG
        \item GraphRAG (static, one hop)
        \item Think-on-Graph (reimplemented)
        \item \textbf{AGR}
      \end{itemize}
      The graph systems start from the datasets' topic entities. The
      control is there because these benchmarks may be partly
      \alert{memorised}.
    \end{column}
    \begin{column}{0.53\textwidth}
      \begin{block}{Held constant}
        \small\agrraggedright
        \textsf{gpt-5.4-mini-2026-03-17} at temperature $0$, the $25$-call
        cap, the questions and the graph.
      \end{block}
      \vspace{1mm}
      \begin{block}{\alert{Not} held equal}
        \small\agrraggedright
        ToG prunes to $40$/$20$ candidates per step, AGR to $300$/$200$.
        \alert{Narrower is cheaper}, so it cannot explain ToG's clipping,
        but its unclipped score is a \alert{lower bound}. The baselines'
        prompt is plain, while AGR's carries grounding rules.
      \end{block}
    \end{column}
  \end{columns}
  \takeaway{Differences are attributable to \textbf{architecture}, not to
    model capacity. Both exceptions are on this slide.}
\end{frame}

% ===================== 17 =====================
\begin{frame}{Main results}
  \bodyshift{-2.75pt}
  \begin{table}
    \small
    % 1.1 rather than the deck's 1.2: the significance sentence under the
    % table is new, and the rows give back the line it needs.
    \renewcommand{\arraystretch}{1.1}
    % The two @{\hspace} separators mark the group boundaries -- WebQSP |
    % CWQ and CWQ | cost -- so the eye pairs the right columns without
    % having to read the spanning headers.
    %
    % "CWQ" rather than "ComplexWebQuestions": as a spanning header the long
    % name was wider than the two columns beneath it and stretched them.
    % "Cost / q" is the same problem in the third group.
    \begin{tabular}{@{}lrr@{\hspace{7mm}}rr@{\hspace{7mm}}rr@{}}
      \toprule
      & \multicolumn{2}{c}{\textbf{WebQSP}}
      & \multicolumn{2}{c}{\textbf{CWQ}}
      & \multicolumn{2}{c}{\textbf{Cost / q}} \\
      \cmidrule(lr){2-3}\cmidrule(lr){4-5}\cmidrule(lr){6-7}
      \textbf{System} & \textbf{Hits@1} & \textbf{F1}
                      & \textbf{Hits@1} & \textbf{F1}
                      & \textbf{Tokens} & \textbf{Calls} \\
      \midrule
      No-retrieval    & 0.453 & 0.305 & 0.307 & 0.279 & 113   & 1.0 \\
      Vector-RAG      & 0.568 & 0.432 & 0.203 & 0.168 & 527   & 1.0 \\
      GraphRAG        & 0.340 & 0.262 & 0.205 & 0.183 & 531   & 1.0 \\
      Think-on-Graph  & 0.660 & 0.477 & 0.435 & 0.378 & 3{,}615 & 12.8 \\
      \textbf{AGR}    & \textbf{0.755} & \textbf{0.642}
                      & \textbf{0.522} & \textbf{0.469}
                      & 4{,}511 & \textbf{6.2} \\
      \bottomrule
    \end{tabular}
  \end{table}
  % Every figure below is bound in check_slides.py: the CWQ costs and the
  % interval widths to thesis_numbers.json, and the McNemar bound to its
  % mcnemar_vs_baselines list. Why no-retrieval outscores GraphRAG on raw
  % hits is said aloud, with its two assertion precisions (section 17 of
  % the script); a fourth line here pushed the takeaway off the slide.
  {\small\agrraggedright Cost is WebQSP. On CWQ, AGR spends $6{,}818$
  tokens and $8.9$ calls, Think-on-Graph $5{,}572$ and $18.2$. All eight
  paired McNemar tests favour AGR ($p < 0.004$), and the 95\% Hits@1
  intervals are $\pm 4$ to $5$ points.\par}
  \takeaway{Ahead of every baseline, at half ToG's calls and 22 to 25\%
    more tokens.}
\end{frame}

% ================ 18 to 25, the research questions ================
% Each slide is titled with its question exactly as sec:rqs words it, on
% one line, and what that slide answers is its subtitle. RQ1 and RQ2 do not
% fit one line at the frame-title size, so each question's slides sit in a
% group that sets its title a little smaller and closes after its last
% slide (preamble.tex says how much, and why a group). RQ3 fits, and its
% two slides have no group.
{\rqtitlesize{14pt}{17.5pt}% RQ1, slides 18 and 19

% ===================== 18 =====================
% The subtitle is the second half of RQ1 as sec:rqs elaborates it, which
% slide 8 carried until the questions there became the book's verbatim.
\begin{frame}{RQ1: Does agentic navigation improve multi-hop factual accuracy?}{And does the advantage grow with hop count?}
  \bodyshift{1.25pt}
  \centering
  \scalebox{0.92}{% !TEX root = ../thesis_defense_0421052099.tex
% GENERATED by scripts/build_figures.py from
% results/phase4/thesis_numbers.json -- do not edit.

\definecolor{agrNoRet}{HTML}{6E6E6E}%
\definecolor{agrVector}{HTML}{E69F00}%
\definecolor{agrGraph}{HTML}{CC79A7}%
\definecolor{agrToG}{HTML}{009E73}%
\definecolor{agrAGR}{HTML}{0072B2}%
\definecolor{agrWQ}{HTML}{0072B2}%
\definecolor{agrCWQ}{HTML}{D55E00}%
\begin{tikzpicture}
\begin{axis}[agrplot, name=hop1, 
  width=52mm, height=34mm,
  title={WebQSP},
  ylabel={Hits@1},
  xtick={0,1,2}, xticklabels={{1 hop\\($n{=}256$)}, {2 hops\\($n{=}128$)},
    {3+ hops\\($n{=}4$)}},
  xticklabel style={align=center},
  xmin=-0.25, xmax=2.25, ymin=0, ymax=0.9,
  legend style={at={(1.11,-0.56)}, anchor=north, legend columns=5,
    column sep=1.8mm},
  legend cell align=left,
  legend image post style={solid},
]
\addplot[dashed, mark=o, mark size=2.6pt, color=agrNoRet, very thick]
  coordinates {(0,0.43) (1,0.52) (2,0.25)};
\addlegendentry{No-retrieval}
\addplot[dashed, mark=square, mark size=2.6pt, color=agrVector, very thick]
  coordinates {(0,0.82) (1,0.12) (2,0.0)};
\addlegendentry{Vector-RAG}
\addplot[dashed, mark=triangle, mark size=2.6pt, color=agrGraph, very thick]
  coordinates {(0,0.3) (1,0.44) (2,0.25)};
\addlegendentry{GraphRAG}
\addplot[dashed, mark=diamond, mark size=2.6pt, color=agrToG, very thick]
  coordinates {(0,0.63) (1,0.78) (2,0.25)};
\addlegendentry{Think-on-Graph}
\addplot[dashed, mark=*, mark size=2.6pt, color=agrAGR, very thick]
  coordinates {(0,0.75) (1,0.84) (2,0.0)};
\addlegendentry{AGR}
\end{axis}
\begin{axis}[agrplot, name=hop2, at={($(hop1.east)+(18mm,0)$)}, anchor=west,
  width=52mm, height=34mm,
  title={ComplexWebQuestions},
  ylabel={Hits@1},
  xtick={0,1,2}, xticklabels={{1 hop\\($n{=}137$)}, {2 hops\\($n{=}211$)},
    {3+ hops\\($n{=}49$)}},
  xticklabel style={align=center},
  xmin=-0.25, xmax=2.25, ymin=0, ymax=0.9,
%
]
\addplot[mark=o, mark size=2.6pt, color=agrNoRet, very thick]
  coordinates {(0,0.38) (1,0.3) (2,0.14)};
\addplot[mark=square, mark size=2.6pt, color=agrVector, very thick]
  coordinates {(0,0.33) (1,0.14) (2,0.1)};
\addplot[mark=triangle, mark size=2.6pt, color=agrGraph, very thick]
  coordinates {(0,0.34) (1,0.16) (2,0.02)};
\addplot[mark=diamond, mark size=2.6pt, color=agrToG, very thick]
  coordinates {(0,0.53) (1,0.37) (2,0.45)};
\addplot[mark=*, mark size=2.6pt, color=agrAGR, very thick]
  coordinates {(0,0.46) (1,0.55) (2,0.57)};
\end{axis}
\end{tikzpicture}%
}

  % The dashed WebQSP panel is one of the boundaries the figure's own
  % caption states, and the slide used to leave it unexplained.
  % Think-on-Graph's one-hop lead on CWQ (0.53 against 0.46) is the other,
  % and it is said aloud.
  {\small On CWQ, AGR goes $0.46 \to 0.55 \to 0.57$, the \textbf{only}
  system ending above where it started. Three of the other four decay
  monotonically. WebQSP is dashed ($n{=}4$ at 3+).\par}
  \takeaway{Yes. The evidence is a shape, not a difference of means.}
\end{frame}

% ===================== 19 =====================
\begin{frame}{RQ1: Does agentic navigation improve multi-hop factual accuracy?}{The caveat I want to raise myself}
  \begin{table}
    \small
    % 1.1, as slides 20, 23 and 24 set it: the question at 14pt cost this
    % slide 2.6pt it did not have.
    \renewcommand{\arraystretch}{1.1}
    % Grouped like slide 17's table, and for the same reason: the long
    % spanning header was wider than the two columns under it.
    \begin{tabular}{@{}lrr@{\hspace{7mm}}rr@{}}
      \toprule
      & \multicolumn{2}{c}{\textbf{WebQSP}}
      & \multicolumn{2}{c}{\textbf{CWQ}} \\
      \cmidrule(lr){2-3}\cmidrule(lr){4-5}
      \textbf{Hits@1 on\ldots} & \textbf{ToG} & \textbf{AGR}
                               & \textbf{ToG} & \textbf{AGR} \\
      \midrule
      questions ToG \emph{finishes} & \textbf{0.852} & 0.788
                                    & \textbf{0.629} & 0.607 \\
      questions ToG is \emph{clipped} on & 0.197 & \textbf{0.675}
                                    & 0.188 & \textbf{0.415} \\
      \midrule
      % The denominators are in the note: a \multicolumn wider than the
      % columns it spans widens the last of them.
      clip rate & \multicolumn{2}{c}{$29.2\%$ ($117$)}
                & \multicolumn{2}{c}{$44.0\%$ ($176$)} \\
      \bottomrule
    \end{tabular}
  \end{table}
  % "Width times depth multiplies ... A plan plus a scored frontier does not
  % multiply" (sec:tog-split) is the mechanism behind half the calls, and
  % the slide used to show the effect without it.
  {\small\agrraggedright Where Think-on-Graph finishes inside the shared
  $25$-call cap, it is \textbf{ahead of AGR on both datasets}. The entire
  aggregate margin comes from the questions it cannot finish, $117$ and
  $176$ of $400$. Beam search pays width times depth in calls, and AGR's
  plan and scored frontier do not multiply.\par}
  \takeaway{The claim is efficiency under a fixed budget, not better
    reasoning per step.}
\end{frame}

}% end of RQ1's title size
{\rqtitlesize{11.6pt}{14.5pt}% RQ2, slides 20 to 23

% ===================== 20 =====================
% The subtitle is the question this slide answers, RQ2 in its narrowest
% form, and the script opens on the same words.
\begin{frame}{RQ2: What does pre-generation verification contribute beyond graph navigation?}{Does anything ungrounded get asserted?}
  \bodyshift{-2.75pt}
  \begin{table}
    \small
    \renewcommand{\arraystretch}{1.1}
    % Count and rate merged in one column, so the column is sized by the
    % data rather than by its header.
    \begin{tabular}{@{}lrr@{}}
      \toprule
      \textbf{System} & \textbf{Asserted} & \textbf{Ungrounded} \\
      \midrule
      No-retrieval    & $1{,}001$ & $221$ ($22.1\%$) \\
      Vector-RAG      & $1{,}329$ & $45$ ($3.4\%$) \\
      GraphRAG        & $1{,}024$ & $8$ ($0.8\%$) \\
      Think-on-Graph  & $1{,}265$ & $\mathbf{0}$ ($\mathbf{0.0\%}$) \\
      \textbf{AGR}    & $1{,}709$ & $\mathbf{0}$ ($\mathbf{0.0\%}$) \\
      \bottomrule
    \end{tabular}
  \end{table}
  % The Tier-2 band is tab:groundedness's own: "zero structural
  % hallucination coexists with roughly a third to a half of assertions
  % failing a strict semantic-support test". Its bounds are held to
  % groundedness_tier2_judge in thesis_numbers.json.
  {\small\agrraggedright Both datasets pooled. A stricter semantic check,
  on a sample of each system's answers, supports only $36.7$ to
  $66.7\%$ of them.\par}
  \takeaway{\alert{Think-on-Graph reaches $0.0\%$ too}. So this is a
    property of \emph{graph navigation}, not of the verification layer.}
\end{frame}

% ===================== 21 =====================
% Three negative findings, each the book's. The middle one is the second
% part of sec:findings' RQ2 answer, which the deck did not carry: the
% precision lead "belongs to the architecture", the layer's removal does not
% move it, and "some share of the precision margin ... may be prompt rather
% than structure". 77/80 is sec:ablations' development figure, and the
% 2/398 is computed from the paired records by check_slides.py.
\begin{frame}{RQ2: What does pre-generation verification contribute beyond graph navigation?}{What verification does \emph{not} do}
  \begin{columns}[T]
    \begin{column}{0.48\textwidth}
      \begin{block}{It does not move accuracy}
        \small\agrraggedright
        Remove it and Hits@1 and F1 do not detectably change, at
        $p = 1.0$ on \emph{both} datasets. It passed $77$ of $80$
        development drafts untouched and changed $2$ of $398$ test answers,
        one each way.
      \end{block}
      \vspace{1mm}
      \begin{block}{Its evidence ends with the run}
        \small\agrraggedright
        The log keeps the \alert{count} of supporting triples and drops the
        list itself.
      \end{block}
    \end{column}
    \begin{column}{0.48\textwidth}
      \begin{block}{It does not explain precision}
        \small\agrraggedright
        On CWQ, AGR is right on $67.9\%$ of what it asserts, Think-on-Graph
        on $56.1\%$. Removing the layer moves AGR's precision by a point at
        most. The lead belongs to the architecture, and possibly in part to
        its drafting prompt.
      \end{block}
    \end{column}
  \end{columns}
  \vspace{2mm}\sinkfoot
  {\small Not a retraction of slide 17. That was AGR against four
  baselines. This is AGR against \emph{itself} with one part removed.
  \alert{What it does buy is the next slide.}\par}
\end{frame}

% ===================== 22 =====================
% The hedge deltas stay, because they are what removing the layer measured.
% What they mean is the paired records', computed in check_slides.py: the
% six CWQ questions the layer withheld are exactly the six the ablated run
% answered, and none of those answers was correct; the one WebQSP question
% it withheld, the ablated run got right. sec:verification-tradeoff
% establishes the layer's value "at case level", and the De Niro case of
% sec:answer-selection is its specimen of a correct rejection.
\begin{frame}{RQ2: What does pre-generation verification contribute beyond graph navigation?}{So what does it do?}
  \bodyshift{-3.25pt}
  \begin{columns}[T]
    \begin{column}{0.48\textwidth}
      \begin{block}{It withholds the ungrounded}
        \small\agrraggedright
        Remove it and hedging falls $23.2\% \to 20.2\%$ on CWQ. Those six
        answers were all wrong. It falls $8.5\% \to 8.0\%$ on WebQSP, one
        answer, which was right.
      \end{block}
      % 0.5mm, not 1mm: the subtitle cost this slide a line, and this
      % column was 0.8pt over.
      \vspace{0.5mm}
      \begin{block}{It attaches evidence}
        \small\agrraggedright
        From \alert{one route of three}, it attaches supporting triples,
        inside the run only. The other two routes,
        \mbox{\textsf{verify\_connection}} and entailment, attach none.
      \end{block}
    \end{column}
    \begin{column}{0.48\textwidth}
      \begin{block}{It rejected a fabricated actor}
        \small\agrraggedright
        A draft said Marlon Brando was in \emph{Joy}. He was not, and the
        verifier rejected the claim. De Niro was, and the evaluator had
        resolved him.
      \end{block}
    \end{column}
  \end{columns}
  \takeaway{Not right more often. Right \alert{for a reason it can show},
    at a small cost in recall.}
\end{frame}

% ===================== 23 =====================
% The between-system evidence that a hedge is not accuracy bought with
% silence, and the answer to the objection the failure census invites
% three slides later: 256 failures, of which 125 are hedges.
%
% The takeaway may not overstate CWQ: AGR hedges 23.0 there against
% Think-on-Graph's 22.5, which is half a point MORE, not less.
%
% The population note is not optional. The previous slide's 8.5 and 23.2
% are the ablation half-samples and these are the full 400 per dataset.
\begin{frame}{RQ2: What does pre-generation verification contribute beyond graph navigation?}{Does it just refuse more often?}
  \bodyshift{-0.5pt}
  \begin{table}
    \small
    % 1.1, as slides 20 and 24 set it: the subtitle cost this slide a line.
    \renewcommand{\arraystretch}{1.1}
    \begin{tabular}{@{}lrr@{}}
      \toprule
      \textbf{System} & \textbf{WebQSP} & \textbf{CWQ} \\
      \midrule
      No-retrieval    & $12.2\%$ & $38.0\%$ \\
      Vector-RAG      & $27.0\%$ & $48.2\%$ \\
      GraphRAG        & $55.8\%$ & $60.8\%$ \\
      Think-on-Graph  & $18.8\%$ & $22.5\%$ \\
      \textbf{AGR}    & $\mathbf{8.2\%}$ & $23.0\%$ \\
      \bottomrule
    \end{tabular}
  \end{table}
  {\small A hedge scores zero on Hits@1, exactly as a wrong answer does.
  AGR hedges \alert{least of the five} on WebQSP, where it is also the most
  accurate. These are the full $400$ questions per dataset. The previous
  slide's $8.5$ and $23.2$ are the ablation half-samples.\par}
  \takeaway{Not more often. The lowest rate of the five on WebQSP, and on
    CWQ within half a point of \mbox{Think-on-Graph} while more accurate.}
\end{frame}

}% end of RQ2's title size: from here on titles are at full size

% ===================== 24 =====================
% The subtitle is what the table is, in the script's own words.
\begin{frame}{RQ3: Which components contribute what, at what token cost?}{Four ablations, each a paired test against the full system}
  \begin{table}
    \small
    % 1.05, not the 1.1 of slides 20 and 23: the question at full size
    % cost this slide 0.95pt it did not have.
    \renewcommand{\arraystretch}{1.05}
    % Three groups, separated like slide 17's: the two datasets and the
    % token column.
    \begin{tabular}{@{}lrr@{\hspace{7mm}}rr@{\hspace{7mm}}r@{}}
      \toprule
      & \multicolumn{2}{c}{\textbf{WebQSP}}
      & \multicolumn{2}{c}{\textbf{CWQ}} & \\
      \cmidrule(lr){2-3}\cmidrule(lr){4-5}
      \textbf{Condition} & \textbf{F1} & \textbf{$p$}
                         & \textbf{F1} & \textbf{$p$}
                         & \textbf{WebQSP tokens} \\
      \midrule
      Full system     & 0.659 & ---   & 0.445 & ---   & 4{,}562 \\
      No planner      & \textbf{0.742} & \alert{\textbf{0.006}}
                      & 0.398 & 0.088 & \textbf{3{,}159} \\
      No backtracking & 0.646 & 0.727 & 0.445 & 1.000 & 4{,}261 \\
      No verifier     & 0.663 & 1.000 & 0.436 & 1.000 & 4{,}460 \\
      Embedding-only scoring & 0.643 & 0.664 & 0.437 & 0.481 & 3{,}413 \\
      \bottomrule
    \end{tabular}
  \end{table}
  \vspace{-2mm}
  % "Full system" is the main run restricted to the same halves, which is
  % why its WebQSP F1 is 0.659 here and 0.642 on slide 17. The p-values are
  % uncorrected and read as descriptive (sec:ablations). The test is on
  % per-question correctness, so "on Hits@1" stops the p from being read as
  % a test of the F1 beside it. "Tokens are for WebQSP" is the column's
  % header now: the subtitle cost this slide a line, and the note's third
  % line held only those two words.
  {\small\agrraggedright $p$ is a paired McNemar test on Hits@1, $n =
  200$/$198$, uncorrected. \emph{Full system} is the main run on the same
  halves, and \emph{embedding-only} drops the model term.\par}
  \takeaway{Three of the four change nothing detectable. The fourth is on
    the next slide, and its sign is backwards.}
\end{frame}

% ===================== 25 =====================
% The strata and the mechanism together. sec:ablations gives the per-stratum
% Hits@1 (held in check_slides.py against ablations.by_hop_stratum), and
% sec:decomposition-errors names context-stripping as the mechanism the
% census "would feature if it could feature only one". WebQTrn-2570 is its
% strongest specimen.
% Full-width blocks, one above the other. Side by side at half the measure
% the first title wrapped "WebQSP" onto a line of its own, which is the
% defect the 0.61 column width of the September deck existed to prevent.
\begin{frame}{RQ3: Which components contribute what, at what token cost?}{One effect, and its sign is backwards}
  \bodyshift{1.75pt}
  \begin{block}{Removing the planner \emph{improves} WebQSP}
    \small\agrraggedright
    $+0.083$ F1, $p = 0.006$, and $-31\%$ tokens. CWQ trends the other
    way, $-0.047$ at $p = 0.088$. Without the planner, WebQSP rises at one
    hop ($0.75 \to 0.85$) and at two ($0.84 \to 0.89$), while CWQ's two-hop
    stratum falls $0.48 \to 0.34$.
  \end{block}
  \vspace{1mm}
  \begin{block}{How decomposition hurts}
    \small\agrraggedright
    It can strip context. ``This 33rd president was the nation's leader
    during WW2'' became two lookups, WW2 was never used, and the run
    answered Woodrow Wilson. Undecomposed, it found Truman in three calls.
  \end{block}
  \takeaway{\textbf{Gate decomposition on structure.} The measurement forced
    that conclusion rather than confirming one.}
\end{frame}

% ===================== 26 =====================
\begin{frame}{Every failure, read and labelled}
  \bodyshift{-1pt}
  \begin{columns}[T]
    \begin{column}{0.44\textwidth}
      % Provenance, from app:census-denominator: AGR fails on 98 of 400
      % WebQSP and 191 of 400 CWQ questions; 33 and 34 of those are set
      % aside as adjudicated benchmark defects or near-misses; a separate
      % reading of the 36 planner-discordant questions is added back, less
      % the two of them that are adjudicated defects. 85 + 171.
      All $256$, read and labelled by hand.

      % The lead-in is its own paragraph and the table starts 2mm under
      % it: at the old spacing the top rule sat 1.4pt below the text and
      % read as a partial underline. The last row closes the column to the
      % heading's 256 -- the six named rows alone sum to 220, and
      % check_slides.py now holds the sum.
      \vspace{2mm}
      \begin{tabular}{@{}lr@{}}
        \toprule
        \textbf{Category} & \textbf{n} \\
        \midrule
        Relation selection   & 64 \\
        Composite claim      & 47 \\
        Knowledge-graph gap  & 43 \\
        Decomposition error  & 38 \\
        Answer selection     & 15 \\
        \alert{Echo attractor} & \alert{13} \\
        \midrule
        Six smaller categories & 36 \\
        \bottomrule
      \end{tabular}
    \end{column}
    \begin{column}{0.52\textwidth}
      % These are the thesis's own Total column -- but sec:taxonomy keeps
      % wrong and hedge apart because "a pooled percentage would describe
      % neither", and pooling hides the shape flip. "Composite claim" rather
      % than \texttt{composite\_claim}: the monospace token cannot
      % hyphenate on this measure.
      \begin{block}{Totals hide the shape}
        \small\agrraggedright
        The $256$ are AGR's failures over the $800$ test questions, less
        the benchmark defects adjudicated before the census and the
        surface-form near-misses. Wrong and hedge are \emph{never pooled}
        in the thesis. The shape flips too. Composite claim is $1$ on
        WebQSP against $46$ on CWQ.
      \end{block}
      \vspace{1mm}
      {\small\agrraggedright The largest category is \alert{relation
      selection}, the agent taking the wrong edge rather than inventing
      one.\par}
    \end{column}
  \end{columns}
  \takeaway{The characteristic error of a graph navigator is \emph{not}
    invention.}
\end{frame}

% ===================== 27 =====================
% This read "it appears across systems, so it is a property of the
% task, not of AGR" -- defensive where the thesis is substantive.
% Real specimens rather than invented ones. WebQTest-1707 is sec:echo's own
% "topic" example. The cross-system half is about evaluation, not blame:
% sec:echo calls the attractor "invisible to any evaluation treating
% systems as independent", and sec:contribution says the contribution is
% the named mechanism "and what it means for consensus-based evaluation
% rather than the frequency".
\begin{frame}{The echo attractor}
  \begin{columns}[T]
    \begin{column}{0.5\textwidth}
      \agrraggedright
      Sixth in that census, with $13$ cases, and the one worth naming. A
      real, \alert{grounded} entity \alert{one hop} from the answer.
      \vspace{2mm}

      {\small Asked who plays Lex Luthor on \emph{Smallville}, AGR named
      John Glover, who plays his father Lionel.\par}
    \end{column}
    \begin{column}{0.46\textwidth}
      \begin{block}{Why it matters}
        \agrraggedright
        Any grounding check \emph{passes} it, since it is true. The entity
        is real, it was traversed, and the triple exists. By construction,
        verification cannot catch it.
      \end{block}
    \end{column}
  \end{columns}
  \vspace{2mm}
  {\agrraggedright Different systems fall into it \alert{together}, so no
  evaluation treating them as independent can see it. A policy of
  rescoring on majority agreement turns it into apparent
  \alert{correctness}.\par}
  \takeaway{The contribution is the mechanism, named, not the count.}
\end{frame}

% ===================== 28 =====================
% sec:findings answers each question in its own words, and the deck never
% said the answers. The takeaway is the chapter's closing sentence, broken
% into three.
\begin{frame}{The three answers}
  \small
  % A margin as wide as the label, as on slide 8, so the labels start at
  % the text block's edge, flush with the takeaway bar. At the default
  % margin they hung 4pt outside both.
  \settowidth{\leftmargini}{\textbf{RQ1}}%
  \addtolength{\leftmargini}{\labelsep}%
  \begin{enumerate}
    \item[\textbf{RQ1}] \textbf{Yes.} AGR leads every accuracy cell, all
      eight paired tests reject, and on CWQ only AGR rises with hop count.
      The margin over Think-on-Graph lies in the budget.
      \vspace{1mm}
    \item[\textbf{RQ2}] Navigation, not verification, removes structural
      hallucination. The precision lead belongs to the architecture. The
      layer buys auditability, not accuracy.
      \vspace{1mm}
    \item[\textbf{RQ3}] Only the planner has a detectable effect, and its
      sign depends on the stratum. The whole system runs at half the
      agentic baseline's calls.
  \end{enumerate}
  \takeaway{Navigation is where the accuracy comes from. Verification is
    where the auditability comes from. Decomposition is sometimes harmful.}
\end{frame}

% ===================== 29 =====================
\begin{frame}{Contributions}
  \bodyshift{-3.25pt}
  \begin{columns}[T]
    \begin{column}{0.50\textwidth}
      % THESE ARE THE THESIS'S SIX, IN ITS ORDER (sec:contribution: the
      % framework paragraph, four \textbf lead-ins, the protocol paragraph).
      % check_slides.py holds each one against sec:contribution. Slide 8
      % previews the same six.
      %
      % "pre-specified": nothing was filed with a registry, and the thesis
      % uses the same word. The protocol item names the decision that was
      % missed, as sec:contribution does ("including the one narrowly
      % missed"); its kappa is held to judge_validation in the JSON.
      \small\let\raggedright\agrraggedright
      \begin{enumerate}\setlength{\itemsep}{0pt}
        \item AGR and its \alert{Structural Verification Layer}
        \item A \alert{component-level ablation} of four mechanisms
        \item \alert{Stratum-dependent decomposition}, gated on structure
        \item The \emph{\alert{echo attractor}}, named
        \item \alert{Benchmark-defect rates} for $57$ questions with wrong
          or \mbox{ambiguous gold}
        \item A \alert{pre-specified} protocol of four decisions, one
          missed ($\kappa = 0.6995$ against $0.7$)
      \end{enumerate}
    \end{column}
    \begin{column}{0.46\textwidth}
      % The first three map to sec:limitations-final's heads, in its
      % severity order, which check_slides.py reads and holds this list to.
      % The fourth is sec:scope's "How to read the title": the registered
      % title is broader than the findings, and the thesis says how to
      % read it. "At face value" is its "better read against the findings
      % than taken as claims in their own right". The item used to list the
      % three terms and never say what was not claimed about them; the list
      % is the script's if-asked note now. With the list in brackets the
      % last line's descenders overran this frame by 1.7pt.
      \begin{block}{What I would not claim}
        \small\let\raggedright\agrraggedright
        \begin{itemize}\setlength{\itemsep}{0pt}
          \item The verifier logs only what it \alert{rejects}. Wrongful
            acceptance is unmeasured, and the evidence is not persisted.
          \item One environment, one backbone, one run
          \item Think-on-Graph leads where it finishes, from a
            \alert{narrower candidate set} ($40$/$20$ against
            $300$/$200$)
          \item The title's terms at face value
        \end{itemize}
      \end{block}
    \end{column}
  \end{columns}
  \centrebody
\end{frame}

% ===================== 30 =====================
% conclusion.tex's future work, the three repairs the census earned and the
% operators the categories specify, and its closing remark as the takeaway.
% The counts are the census's own and check_slides.py holds them.
% Full width, so each item is one line: in two columns, every
% parenthetical count wrapped away from the noun it counts. Two headed
% lists rather than two blocks, because the block chrome was the 15pt the
% takeaway needed.
\begin{frame}{What comes next}
  \bodyshift{4.5pt}
  \small\let\raggedright\agrraggedright
  \textbf{Repairs the census earned}
  \begin{itemize}\setlength{\itemsep}{0pt}\setlength{\parskip}{0pt}
    \item Tell the decomposer which argument is the answer ($9$
      extraction-bug cases)
    \item Log accepted claims with the triples that matched them
    \item Re-run the baselines at equal candidate widths and a wider radius
  \end{itemize}
  \textbf{Operators the failures specify}
  \begin{itemize}\setlength{\itemsep}{0pt}\setlength{\parskip}{0pt}
    \item Gate decomposition on question structure
    \item A set-intersection operator ($47$ composite-claim cases)
    \item Ordinals and literal values ($43$ knowledge-graph-gap cases)
  \end{itemize}
  \takeaway{Fact verification against a knowledge graph is the easier half.
    \alert{Relevance} verification is the half still open.}
\end{frame}

% ===================== CLOSE =====================
{\setbeamertemplate{footline}{}
\begin{frame}[plain]
  \vfill
  \centering
  {\Large\bfseries\color{agrdark} Thank you}\\[4mm]
  {\large Questions}
  \vfill
\end{frame}}

% ===================== BACKUP =====================
% Not presented. These follow the closing slide so a question can be
% answered by paging past it, rather than by switching to a second PDF.
% A backup frame with no takeaway bar closes with \centrebody, as the
% presented ones do. Without it the frametitle's glue is the only stretch
% in the frame and the body sinks to the floor: "which budgets actually
% bind" sat 82pt under its title and 2pt above the page number.
\begin{frame}{Backup: budget configuration}
  \bodyshift{-4.5pt}
  \begin{table}
    \small
    \begin{tabular}{@{}lrL{52mm}@{}}
      \toprule
      \textbf{Parameter} & \textbf{Value} & \textbf{Enforced in} \\
      \midrule
      \texttt{max\_depth}        & 4 & both routers \\
      \texttt{beam\_width}       & 3 & explorer (slice) \\
      \texttt{max\_anchors}      & 5 & explorer \emph{and} evaluator \\
      \texttt{max\_backtracks}   & 3 & post-evaluation router \\
      \texttt{max\_llm\_calls}   & 25 & LLM client \\
      \texttt{max\_verify\_iters}& 2 & post-verification router, verifier \\
      \texttt{max\_seconds}      & 300 & both routers \\
      \bottomrule
    \end{tabular}
  \end{table}
  {\small Termination does not depend on model behaviour. Every cycle passes
  through a router that checks a monotone counter.\par}
  \centrebody
\end{frame}

\begin{frame}{Backup: which budgets actually bind}
  \bodyshift{-3.5pt}
  \begin{table}
    \small
    \begin{tabular}{@{}lrrr@{}}
      \toprule
      \textbf{Budget} & \textbf{WebQSP} & \textbf{CWQ} & \textbf{Pooled} \\
      \midrule
      Depth cap        & $16.5\%$ & $34.8\%$ & $25.6\%$ \\
      Backtrack cap    & $3.2\%$  & $9.2\%$  & $6.2\%$ \\
      Verify-iteration cap & $1.0\%$ & $2.8\%$ & $1.9\%$ \\
      \textbf{Call cap} & $\mathbf{0.0\%}$ & $\mathbf{0.0\%}$
                        & $\mathbf{0.0\%}$ \\
      \bottomrule
    \end{tabular}
  \end{table}
  {\small Percentage of questions on which each budget refused a request.
  AGR never reaches the $25$-call cap. That is what makes a clipped
  Think-on-Graph a meaningful comparison.\par}
  \centrebody
\end{frame}

\begin{frame}{Backup: the full failure census}
  \bodyshift{3pt}
  \centering
  \resizebox{\textwidth}{!}{%
    % !TEX root = ../thesis_defense_0421052099.tex
% GENERATED by scripts/build_figures.py from
% results/phase4/thesis_numbers.json -- do not edit.

\definecolor{agrNoRet}{HTML}{6E6E6E}
\definecolor{agrVector}{HTML}{E69F00}
\definecolor{agrGraph}{HTML}{CC79A7}
\definecolor{agrToG}{HTML}{009E73}
\definecolor{agrAGR}{HTML}{0072B2}
\definecolor{agrWQ}{HTML}{0072B2}
\definecolor{agrCWQ}{HTML}{D55E00}
\begin{tikzpicture}
\begin{axis}[agrplot,
  width=112mm, height=62mm,
  xbar stacked, bar width=3.0mm,
  xlabel={Questions}, xmin=0, xmax=70,
  ytick={0,1,2,3,4,5,6,7,8,9,10,11},
  yticklabels={Relation selection, Composite claim, Knowledge-graph gap,
    Decomposition error, Answer selection, Echo attractor, Ambiguous question,
    Premature termination, Gold noise, Other, Wrongly rejected claim,
    Wrongly accepted claim},
  yticklabel style={font=\small},
  y dir=reverse, ymin=-0.7, ymax=11.7,
  legend style={at={(0.98,0.05)}, anchor=south east},
  legend cell align=left, reverse legend=false,
]
\addplot[fill=agrWQ, draw=agrWQ!80]
  coordinates {(19,0) (1,1) (8,2) (15,3) (1,4) (5,5) (1,6) (1,7) (1,8) (0,9)
    (0,10) (0,11)};
\addlegendentry{WebQSP wrong}
\addplot[fill=agrWQ!35, draw=agrWQ!80]
  coordinates {(6,0) (0,1) (4,2) (11,3) (1,4) (2,5) (1,6) (4,7) (0,8) (1,9)
    (3,10) (0,11)};
\addlegendentry{WebQSP hedge}
\addplot[fill=agrCWQ, draw=agrCWQ!80, postaction={pattern=north east lines,
  pattern color=white}]
  coordinates {(20,0) (23,1) (7,2) (10,3) (8,4) (4,5) (3,6) (1,7) (2,8) (0,9)
    (1,10) (0,11)};
\addlegendentry{CWQ wrong}
\addplot[fill=agrCWQ!35, draw=agrCWQ!80]
  coordinates {(19,0) (23,1) (24,2) (2,3) (5,4) (2,5) (5,6) (2,7) (3,8) (5,9)
    (1,10) (1,11)};
\addlegendentry{CWQ hedge}
\end{axis}
\end{tikzpicture}
}
  \centrebody
\end{frame}

% Token counts tick as 100, 1,000 and 10,000 here, not as powers of ten.
% As 10^2, 10^3 and 10^4 the one digit that tells the ticks apart is the
% exponent, which printed at 6pt, the smallest text in the deck. Set on
% this slide rather than in build_figures.py, so the book's and the
% paper's copies of the figure are untouched.
\begin{frame}{Backup: accuracy against cost}
  \centering
  {\pgfplotsset{every axis/.append style={log ticks with fixed point}}%
   % !TEX root = ../thesis_defense_0421052099.tex
% GENERATED by scripts/build_figures.py from
% results/phase4/thesis_numbers.json -- do not edit.

\definecolor{agrNoRet}{HTML}{6E6E6E}%
\definecolor{agrVector}{HTML}{E69F00}%
\definecolor{agrGraph}{HTML}{CC79A7}%
\definecolor{agrToG}{HTML}{009E73}%
\definecolor{agrAGR}{HTML}{0072B2}%
\definecolor{agrWQ}{HTML}{0072B2}%
\definecolor{agrCWQ}{HTML}{D55E00}%
\begin{tikzpicture}
\begin{axis}[agrplot, name=plot1, 
  width=52mm, height=34mm,
  title={WebQSP},
  xmode=log, log basis x=10,
  xlabel={Mean tokens per question}, ylabel={Hits@1},
  xmin=80, xmax=12000, ymin=0.15, ymax=0.85,
  legend style={at={(1.13,-0.70)}, anchor=north, legend columns=5,
    column sep=1.8mm},
  legend cell align=left,
  legend image post style={solid},
]
\addplot[only marks, mark=o, mark size=3.0pt, color=agrNoRet, very thick]
  coordinates {(113,0.453)};
\addlegendentry{No-retrieval}
\addplot[only marks, mark=square, mark size=3.0pt, color=agrVector, very thick]
  coordinates {(527,0.568)};
\addlegendentry{Vector-RAG}
\addplot[only marks, mark=triangle, mark size=3.0pt, color=agrGraph, very thick]
  coordinates {(531,0.34)};
\addlegendentry{GraphRAG}
\addplot[only marks, mark=diamond, mark size=3.0pt, color=agrToG, very thick]
  coordinates {(3615,0.66)};
\addlegendentry{Think-on-Graph}
\addplot[only marks, mark=*, mark size=3.0pt, color=agrAGR, very thick]
  coordinates {(4511,0.755)};
\addlegendentry{AGR}
\end{axis}
\begin{axis}[agrplot, name=plot2, at={($(plot1.east)+(18mm,0)$)}, anchor=west,
  width=52mm, height=34mm,
  title={ComplexWebQuestions},
  xmode=log, log basis x=10,
  xlabel={Mean tokens per question}, ylabel={Hits@1},
  xmin=80, xmax=12000, ymin=0.15, ymax=0.85,
%
]
\addplot[only marks, mark=o, mark size=3.0pt, color=agrNoRet, very thick]
  coordinates {(118,0.307)};
\addplot[only marks, mark=square, mark size=3.0pt, color=agrVector, very thick]
  coordinates {(535,0.203)};
\addplot[only marks, mark=triangle, mark size=3.0pt, color=agrGraph, very thick]
  coordinates {(470,0.205)};
\addplot[only marks, mark=diamond, mark size=3.0pt, color=agrToG, very thick]
  coordinates {(5572,0.435)};
\addplot[only marks, mark=*, mark size=3.0pt, color=agrAGR, very thick]
  coordinates {(6818,0.522)};
\end{axis}
\end{tikzpicture}%
}

  \vspace{1mm}
  {\small On \emph{tokens} Think-on-Graph is \textbf{not} dominated. It buys
  lower accuracy at a lower token price. On \emph{calls}, which the budget
  meters, it is.\par}
  \takeaway{The frontier depends on which cost you charge.}
\end{frame}

\begin{frame}{Backup: the benchmark was wrong $57$ times}
  \begin{columns}[T]
    \begin{column}{0.5\textwidth}
      \small
      % sec:benchmark-defects makes the detector and the attractor one
      % finding: "The gap between those figures is not measurement noise.
      % It is the echo attractor of sec:echo, operating across systems."
      % 105 is 58 + 47 flagged, 41 is 22 + 19 adjudicated.
      %
      % The breakdown is app:defects': two DISJOINT sets, 41 + 16 = 57. It
      % read "17 inside the census, 1 counted in both" until 2026-10, the
      % accounting from before WebQTrn-64 was promoted to an exclusion.
      % check_slides.py holds all three lines to benchmark_defects.
      \agrraggedright The same cross-system agreement is also a
      \alert{detector}. Consensus flagged $105$ questions and adjudication
      confirmed $41$. The gap is the attractor above, not label noise.
      \begin{itemize}
        \item $41$ excluded before the census
        \item $16$ more found inside it, as gold-noise or ambiguous rows
        \item No question counted in both
        \item[$\Rightarrow$] \textbf{$57$ distinct questions}
      \end{itemize}
    \end{column}
    \begin{column}{0.48\textwidth}
      \begin{block}{Why report it}
        \small\agrraggedright
        Reported defect rates for these two standard benchmarks are rare and
        usually anecdotal. This is a counted rate with a documented
        adjudication procedure and per-question provenance.
      \end{block}
    \end{column}
  \end{columns}
  \takeaway{A contribution that outlives this particular system.}
\end{frame}

% The board asked for this at the pre-defense: the environment is built
% from RoG's published subgraphs, so a direct comparison is owed. It is a
% backup slide because the honest answer is that RoG leads and the
% comparison does not settle an architectural question. Figures are tab:rog
% (sec:rog-comparison). The training-split sizes are evaluation.tex's, and
% check_slides.py holds them there: this slide said 2,830 and 16,900 until
% 2026-10.
\begin{frame}{Backup: AGR against RoG}
  \bodyshift{-2.75pt}
  \begin{table}
    \small
    \begin{tabular}{@{}lrr@{\hspace{7mm}}rr@{}}
      \toprule
      & \multicolumn{2}{c}{\textbf{WebQSP}}
      & \multicolumn{2}{c}{\textbf{CWQ}} \\
      \cmidrule(lr){2-3}\cmidrule(lr){4-5}
      \textbf{System} & \textbf{Hits@1} & \textbf{F1}
                      & \textbf{Hits@1} & \textbf{F1} \\
      \midrule
      RoG (published) & $85.7$ & $70.8$ & $62.6$ & $56.2$ \\
      \textbf{AGR} (this work) & $75.5$ & $64.2$ & $52.2$ & $46.9$ \\
      \bottomrule
    \end{tabular}
  \end{table}
  {\small\agrraggedright \alert{RoG leads on all four}, and the gap is
  outside AGR's confidence intervals. But RoG is \alert{fine-tuned} on the
  training splits of both benchmarks, $2{,}826$ and $27{,}639$ questions,
  while AGR trains on \alert{nothing}. RoG reports over the full test
  splits, AGR over the certified $400$-question samples of those
  splits.\par}
  \takeaway{Same data, different \textbf{class} of system, supervised
    against zero-shot. Worth stating, but it is not an architecture
    comparison, and no protocol here could make it one.}
\end{frame}

% Four departures from the approved proposal, each recorded in the thesis
% where it applies: sec:mcts-clarification, sec:factuality-related,
% sec:metrics-not-run and sec:pre-registration.
\begin{frame}{Backup: what changed from the proposal}
  \begin{table}
    \small
    \begin{tabular}{@{}L{34mm}L{24mm}L{72mm}@{}}
      \toprule
      \textbf{Proposal} & \textbf{Now} & \textbf{Why} \\
      \midrule
      ``MCTS-inspired'' search & Withdrawn
        & No rollouts and no value backpropagation. It is guided best-first
          search with a bounded beam and backtracking. \\
      \midrule
      FactBench as graph and judge & Dropped
        & An evaluator whose facts sit inside the retrieval graph cannot
          measure hallucination. \\
      \midrule
      Path fidelity metric & Deferred
        & Needs gold SPARQL relation chains, which the RoG distribution
          does not carry. \\
      \midrule
      A registered protocol & Pre-specified
        & Four decisions were fixed in writing, but nothing was filed with
          a registry. \\
      \bottomrule
    \end{tabular}
  \end{table}
  \centrebody
\end{frame}
